\section{Сжатое представление леса разбора}
\label{sec:SPPF}
\tikzsetfigurename{CFG_SPPF_}

Как было показано в разделе~\ref{sec:DerivTree}, дерево вывода представляет ровно один вывод цепочки в грамматике.
Если грамматика неоднозначна, то для одной и той же цепочки может существовать несколько деревьев вывода, которые в совокупности называются \emph{лесом разбора}.
Наивное хранение леса разбора как набора отдельных деревьев приводит к экспоненциальному дублированию общих поддеревьев%
\sidenote{
Предлагаем оценить количество деревьев для цепочки $(a \ b)^+$ в грамматике $S \to S \ S \mid a \ S \ b  \mid a \ b$ в зависимости от длины цепочки.
Отдельно заметим, что для грамматики $S \to S \ S \mid a \ S \ b \mid \varepsilon$ количество деревьев бесконечно из-за возможности превратить любую последовательность из $S$ в пустую строку.
}%
.

Сжатое представление леса разбора (Shared Packed Parse Forest, SPPF)~\sidecite{SPPF}, предложенное Рекерсом, позволяет представить весь лес разбора в виде единого ориентированного графа, в котором одинаковые поддеревья не дублируются, а альтернативные выводы <<упаковываются>> под общим узлом.

\begin{definition}[Простое сжатое представление леса разбора (по Рекерсу)]
    \label{def:basicSPPF}
    Будем считать, что грамматика $G = \langle \Sigma, N, P, S \rangle$ задана в BNF-форме: каждое правило имеет вид $X_k \to \alpha_k$ (без использования <<$|$>> для альтернатив в правой части), пронумеровано $k \in \{1, \dots, |P|\}$.
    \emph{Сжатым представлением леса разбора} (SPPF) цепочки $\omega$ длины $n$ называется ориентированный ациклический граф, вершины которого делятся на два типа.
        \begin{enumerate}
            \item \emph{Символьные вершины}. Метка имеет вид $X_{i,j}$, где $X \in \Sigma \cup N \cup \{\varepsilon\}$, а $i, j$~--- позиции во входной цепочке ($0 \leq i \leq j \leq n$). Символьная вершина $X_{i,j}$ обозначает, что из грамматического символа $X$ выводится подстрока $\omega[i..j-1]$.
            \item \emph{Вершины-продукции}. Метка~--- номер правила $k \in \{1, \dots, |P|\}$. Вершина-продукция с номером $k$ соответствует применению $k$-го правила $X_k \to \alpha_k$. Её непосредственный предок~--- символьная вершина $(X_k)_{i,j}$, а непосредственные потомки~--- символьные вершины, соответствующие элементам правой части $\alpha_k$ и покрывающие (возможно, с разбиением) подстроку с $i$ по $j$.
            \item Одна и та же символьная вершина $X_{i,j}$ может иметь несколько вершин-продукций в качестве потомков (\emph{упаковка} альтернатив, от англ. \emph{packing}).
            \item Корнем SPPF является символьная вершина $S_{0,n}$.
            \item Листьями являются символьные вершины вида $a_{i,i+1}$ для терминалов $a \in \Sigma$ и $\varepsilon_i$ для $\varepsilon$.
        \end{enumerate}
\end{definition}

Если два различных вывода порождают нетерминал $X$ на одной и той же подстроке $\omega[i..j-1]$, то соответствующие им символьные вершины совпадают: это одна и та же вершина $X_{i,j}$, к которой <<подвешиваются>> вершины-продукции для каждого из альтернативных выводов.
Таким образом, SPPF является корректным представлением леса разбора: любое дерево вывода может быть восстановлено выбором ровно одной вершины-продукции для каждой символьной вершины, имеющей упаковку.

На практике компактность SPPF обеспечивается двумя механизмами разделения узлов.
\begin{itemize}
    \item \textbf{Разделение вершин-продукций.} При создании вершины-продукции проверяется, существует ли уже вершина-продукция с тем же номером правила и теми же дочерними символьными вершинами. Если да, она переиспользуется.
    \item \textbf{Разделение символьных вершин.} При создании символьной вершины проверяется, существует ли уже символьная вершина с тем же символом, покрывающая ту же часть входа. Если да, она переиспользуется.
\end{itemize}

Рассмотрим несколько примеров SPPF.

\begin{example}
    \label{exmpl:basic_sppf_unambiguous}
    Начнём с однозначной грамматики для языка Дика на одном типе скобок:
    \[
    G_0 = \langle \{a,b\}, \{S\}, S, \{(1)\ S \to a \ S \ b \ S, \ (2)\ S \to \varepsilon\} \rangle.
    \]
    Построим SPPF для цепочки $_0 a_1 b_2 a_3 b_4 a_5 b_6$.
    Поскольку грамматика однозначна, для данной цепочки существует ровно одно дерево вывода, и SPPF вырождается в обычное аннотированное дерево (с той лишь разницей, что узлы несут дополнительную информацию о позициях и номерах продукций).
    Полученный SPPF представлен на рисунке~\ref{fig:sppf_unambiguous}.
	Если удалить из него узлы, несущие информацию о продукциях, то получится в точности классическое дерево вывода из примера~\ref{exmpl:derivation_tree}.

    \begin{figure}
        \begin{center}
            \resizebox{0.9\textwidth}{!}{\begin{tikzpicture}[shorten >=1pt,node distance=1.2cm]
   \node[symbol_node] (s_0_6)   {$(0,S,6)$};
   \node[prod_node] (p_0_1) [below=of s_0_6] {$0$};
   \node[prod_node,draw=none] (dummy1) [below =of p_0_1] {};
   \node[symbol_node] (s_1_1) [below left=of p_0_1] {$(1,S,1)$};
   \node[symbol_node] (s_2_6) [below right=of p_0_1]  {$(2,S,6)$};

   \node[prod_node] (p_0_2) [below right=of s_2_6] {$0$};

   \node[symbol_node] (s_3_3) [below left=of p_0_2] {$(3,S,3)$};
   \node[symbol_node] (s_4_6) [below right=of p_0_2] {$(4,S,6)$};

   \node[prod_node] (p_0_3) [below right =of s_4_6] {0};

   \node[symbol_node] (s_5_5) [below  =of p_0_3] {$(5,S,5)$};
   \node[symbol_node] (s_6_6) [below right=of p_0_3] {$(6,S,6)$};

%   \node[state,draw=none] (dummy1) [below =of s_0_6] {};
%   \node[state,draw=none] (dummy2) [below =of dummy1] {};

%   \node[state,draw=none] (dummy3) [below left=of p_0_3] {};
%   \node[state,draw=none] (dummy4) [below right=of p_0_4] {};


%   \node[symbol_node] (s_0_2) [left=of dummy3] {$(0,S,2)$};
%
%   \node[symbol_node] (s_2_4) [between=s_0_2 and s_4_6] {$(2,S,4)$};


   \node[prod_node] (p_1_1) [below left =of s_1_1] {$1$};
   \node[prod_node] (p_1_2) [below left =of s_3_3] {$1$};
   \node[prod_node] (p_1_3) [below  =of s_5_5] {$1$};
   \node[prod_node] (p_1_4) [below right=of s_6_6] {$1$};



%   \node[prod_node] (p_2_1) [below =of s_1_1] {$2$};
%   \node[prod_node] (p_2_2) [below =of s_3_3] {$2$};
%   \node[prod_node] (p_2_3) [below =of s_5_5] {$2$};

   \node[symbol_node] (eps_6_6) [below right=of p_1_4] {$(6,\varepsilon,6)$};
   \node[symbol_node] (b_5_6)   [left=of eps_6_6] {$(5,b,6)$};
   \node[symbol_node] (eps_5_5) [left =of b_5_6] {$(5,\varepsilon,5)$};
   \node[symbol_node] (a_4_5)   [left=of eps_5_5] {$(4,a,5)$};
   \node[symbol_node] (b_3_4)   [left=of a_4_5] {$(3,b,4)$};
   \node[symbol_node] (eps_3_3) [left =of b_3_4] {$(3,\varepsilon,3)$};
   \node[symbol_node] (a_2_3)   [left=of eps_3_3] {$(2,a,3)$};
   \node[symbol_node] (b_1_2)   [left=of a_2_3] {$(1,b,2)$};
   \node[symbol_node] (eps_1_1) [left =of b_1_2] {$(1,\varepsilon,1)$};
   \node[symbol_node] (a_0_1) [left=of eps_1_1] {$(0,a,1)$};


    \path[->]
    (s_0_6) edge (p_0_1)

    (p_0_1) edge [bend right] (a_0_1)
    (p_0_1) edge (s_1_1)
    (p_0_1) edge  (b_1_2)
    (p_0_1) edge (s_2_6)

    (s_2_6) edge (p_0_2)

    (p_0_2) edge [bend right] (a_2_3)
    (p_0_2) edge (s_3_3)
    (p_0_2) edge (b_3_4)
    (p_0_2) edge (s_4_6)

    (s_4_6) edge (p_0_3)

    (p_0_3) edge (a_4_5)
    (p_0_3) edge (s_5_5)
    (p_0_3) edge (b_5_6)
    (p_0_3) edge (s_6_6)

    (s_1_1) edge (p_1_1)
    (p_1_1) edge (eps_1_1)

    (s_3_3) edge (p_1_2)
    (p_1_2) edge (eps_3_3)

    (s_5_5) edge (p_1_3)
    (p_1_3) edge (eps_5_5)

    (s_6_6) edge (p_1_4)
    (p_1_4) edge (eps_6_6)


    ;
\end{tikzpicture}
}
        \end{center}
        \caption{SPPF для цепочки $ababab$ в однозначной грамматике $G_0$}
        \label{fig:sppf_unambiguous}
    \end{figure}
\end{example}

\begin{example}
    \label{exmpl:basic_sppf_ambiguous}
    Теперь рассмотрим неоднозначную грамматику, задающую тот же язык правильных скобочных последовательностей (см. раздел~\ref{chpt:CFG}), с пронумерованными продукциями:
    \[
    G_1 = \langle \{a,b\}, \{S\}, S, \{(1)\ S \to S \ S, \ (2)\ S \to a \ S \ b, \ (3)\ S \to \varepsilon\} \rangle.
    \]
    Для цепочки $ababab$ существуют два различных левосторонних вывода (первым применяется продукция $(1)$, после чего первый образовавшийся нетерминал $S$ может быть переписан либо по $(1)$, либо по $(2)$):

    \begin{align*}
        \mathbf{S} \xrightarrow{(1)} \mathbf{S}S \xrightarrow{(1)} \mathbf{S}SS \xrightarrow{(2)} a\mathbf{S}bSS \xrightarrow{(3)} ab\mathbf{S}S \xrightarrow{(2)} aba\mathbf{S}bS \xrightarrow{(3)} abab\mathbf{S} \xrightarrow{(2)} ababa\mathbf{S}b \xrightarrow{(3)} ababab
        \\
        \mathbf{S} \xrightarrow{(1)} \mathbf{S}S \xrightarrow{(2)} a\mathbf{S}bS \xrightarrow{(3)} ab\mathbf{S} \xrightarrow{(1)} ab\mathbf{S}S \xrightarrow{(2)} aba\mathbf{S}bS \xrightarrow{(3)} abab\mathbf{S} \xrightarrow{(2)} ababa\mathbf{S}b \xrightarrow{(3)} ababab
    \end{align*}

    Аннотированные деревья, соответствующие этим двум выводам, показаны на рисунках~\ref{fig:sppf_tree1} и~\ref{fig:sppf_tree2}.

    \begin{figure}
        \begin{center}
            \resizebox{0.9\textwidth}{!}{\input{part_02_Foundations/chapter_06_ContextFreeLanguages/figures/05_SPPF/05_SPPF_fig_02.tex}}
        \end{center}
        \caption{Дерево вывода для первого левостороннего вывода}
        \label{fig:sppf_tree1}
    \end{figure}

    \begin{figure}
        \begin{center}
            \resizebox{0.9\textwidth}{!}{\begin{tikzpicture}[shorten >=1pt,on grid,auto,node distance=1.8cm]
   \node[symbol_node] (s_0_6)   {$S_{0,6}$};
   \node[prod_node,draw=none] (p_0_1) [below left=of s_0_6] {};
   \node[prod_node] (p_0_2) [below right=of s_0_6] {$(1)$};
   \node[symbol_node,draw=none] (s_0_4) [below left=of p_0_1]  {};
   \node[symbol_node] (s_2_6) [below right=of p_0_2]  {$S_{2,6}$};
   \node[prod_node,draw=none] (p_0_3) [below =of s_0_4] {};
   \node[prod_node] (p_0_4) [below =of s_2_6] {$(1)$};

   \node[state,draw=none] (dummy1) [below =of s_0_6] {};
   \node[state,draw=none] (dummy2) [below =of dummy1] {};

   \node[state,draw=none] (dummy3) [below left=of p_0_3] {};
   \node[state,draw=none] (dummy4) [below right=of p_0_4] {};


   \node[symbol_node] (s_0_2) [left=of dummy3] {$S_{0,2}$};
   \node[symbol_node] (s_4_6) [right=of dummy4] {$S_{4,6}$};
   \node[symbol_node] (s_2_4) [between=s_0_2 and s_4_6] {$S_{2,4}$};

   \node[prod_node] (p_1_1) [below =of s_0_2] {$(2)$};
   \node[prod_node] (p_1_2) [below =of s_2_4] {$(2)$};
   \node[prod_node] (p_1_3) [below =of s_4_6] {$(2)$};

   \node[symbol_node] (s_1_1) [below =of p_1_1] {$S_{1,1}$};
   \node[symbol_node] (s_3_3) [below =of p_1_2] {$S_{3,3}$};
   \node[symbol_node] (s_5_5) [below =of p_1_3] {$S_{5,5}$};

   \node[prod_node] (p_2_1) [below =of s_1_1] {$(3)$};
   \node[prod_node] (p_2_2) [below =of s_3_3] {$(3)$};
   \node[prod_node] (p_2_3) [below =of s_5_5] {$(3)$};

   \node[symbol_node] (eps_1_1) [below =of p_2_1] {$\varepsilon_{1}$};
   \node[symbol_node] (eps_3_3) [below =of p_2_2] {$\varepsilon_{3}$};
   \node[symbol_node] (eps_5_5) [below =of p_2_3] {$\varepsilon_{5}$};

   \node[symbol_node] (a_0_1) [left=of eps_1_1] {$a_{0,1}$};
   \node[symbol_node] (a_2_3) [left=of eps_3_3] {$a_{2,3}$};
   \node[symbol_node] (a_4_5) [left=of eps_5_5] {$a_{4,5}$};

   \node[symbol_node] (b_1_2) [right=of eps_1_1] {$b_{1,2}$};
   \node[symbol_node] (b_3_4) [right=of eps_3_3] {$b_{3,4}$};
   \node[symbol_node] (b_5_6) [right=of eps_5_5] {$b_{5,6}$};


    \path[->]
    (s_0_6) edge (p_0_2)
    (p_0_2) edge (s_0_2)
    (p_0_2) edge (s_2_6)
    (s_2_6) edge (p_0_4)
    (p_0_4) edge (s_2_4)
    (p_0_4) edge (s_4_6)

    (s_0_2) edge (p_1_1)
    (s_2_4) edge (p_1_2)
    (s_4_6) edge (p_1_3)

    (p_1_1) edge [bend right] (a_0_1)
    (p_1_1) edge (s_1_1)
    (p_1_1) edge [bend left] (b_1_2)

    (p_1_2) edge [bend right] (a_2_3)
    (p_1_2) edge (s_3_3)
    (p_1_2) edge [bend left] (b_3_4)

    (p_1_3) edge [bend right] (a_4_5)
    (p_1_3) edge (s_5_5)
    (p_1_3) edge [bend left] (b_5_6)

    (s_1_1) edge (p_2_1)
    (p_2_1) edge (eps_1_1)

    (s_3_3) edge (p_2_2)
    (p_2_2) edge (eps_3_3)

    (s_5_5) edge (p_2_3)
    (p_2_3) edge (eps_5_5)

    ;
\end{tikzpicture}
}
        \end{center}
        \caption{Дерево вывода для второго левостороннего вывода}
        \label{fig:sppf_tree2}
    \end{figure}

    В построенных деревьях большое количество символьных вершин совпадает.
    Например, вершины $S_{1,1}$, $S_{3,3}$, $S_{5,5}$ и все терминальные листья в точности одинаковы в обоих деревьях.
    Объединив их в соответствии с определением~\ref{def:basicSPPF}, получим единую структуру, представленную на рисунке~\ref{fig:sppf_ambiguous}.

    \begin{figure}
        \begin{center}
            \resizebox{0.9\textwidth}{!}{\input{part_02_Foundations/chapter_06_ContextFreeLanguages/figures/05_SPPF/05_SPPF_fig_04.tex}}
        \end{center}
        \caption{SPPF для цепочки $ababab$ в неоднозначной грамматике $G_1$}
        \label{fig:sppf_ambiguous}
    \end{figure}
\end{example}

\begin{example}
    \label{exmpl:basic_sppf_cycles}
    Вернёмся к неоднозначной грамматике $G_1$ и рассмотрим более короткую цепочку $ab$ (позиции $0 \rng 2$).
    Данная цепочка может быть выведена через продукцию $(2)$ как $S \to a\ S\ b$, где внутренний $S$ выводит $\varepsilon$ по правилу $(3)$, либо через продукцию $(1)$ как $S \to S\ S$, где левый $S$ выводит $\varepsilon$ по $(3)$, а правый~--- $ab$ через $(2)$ и $(3)$.
    Кроме того, правило $(1)$ в комбинации с $\varepsilon$-выводами позволяет породить бесконечно много выводов через $S \derives[+] S$: например, $S \xrightarrow{(1)} S S \xrightarrow{(3)} \varepsilon S \xrightarrow{(1)} \varepsilon S S \xrightarrow{(3)} \cdots \xrightarrow{(2)} ab$.

    В SPPF это проявляется как наличие циклов: символьная вершина $S_{0,2}$ может иметь вершину-продукцию $(1)$ с дочерними вершинами $S_{0,0}$ и $S_{0,2}$ (либо $S_{0,2}$ и $S_{2,2}$), что зацикливает граф.
    Таким образом, SPPF в общем случае не является деревом и может содержать циклы, вызванные $\varepsilon$-правилами.
    Более подробно циклические структуры в SPPF для этой же грамматики и цепочки будут рассмотрены ниже на примере RSM-варианта SPPF.
\end{example}

Определение~\ref{def:basicSPPF}, восходящее к Рекерсу~\sidecite{SPPF}, даёт интуитивно понятное представление леса разбора через символьные вершины и вершины-продукции.
В дальнейшем Билло и Ланг~\sidecite{Billot:1989:SSF:981623.981641} предложили \emph{бинаризованный SPPF}, гарантирующий объём $O(n^3)$, где $n$~--- длина входной цепочки.
Идея состоит в том, чтобы каждое правило с правой частью длины $k \geq 2$ рассматривалось как последовательность $k-1$ бинарных склеек: для каждой пары соседних подцепочек вводится промежуточный узел, хранящий позицию разделения.
Благодаря этому число возможных узлов-диапазонов ограничено $O(n^2)$, а каждый узел имеет не больше двух детей, что и даёт кубическую оценку общего размера SPPF.

Впоследствии бинаризованный вариант SPPF нашёл применение в алгоритмах обобщённого синтаксического анализа: RNGLR~\sidecite{Scott:2006:RNG:1146809.1146810}, BRNGLR~\sidecite{Scott:2007:BCT:1289813.1289815} и GLL~\sidecite{Scott:2010:GP:1860132.1860320}.
Для работы с EBNF-грамматиками без преобразования в BNF Скотт и Джонстон разработали вариант SPPF, описанный в~\sidecite{SCOTT2018120}.

В определении~\ref{def:basicSPPF} активно используется тот факт, что грамматика находится в BNF. Без этого идея с нумерацией продукций не сработает.
Но мы хотим уметь работать с грамматиками в EBNF без их преобразования. Поэтому далее мы будем использовать вариант SPPF, непосредственно связанный с рекурсивным автоматом из раздела~\ref{sec:RecursiveAutomata}: он наиболее удобен для алгоритмов синтаксического анализа, работающих с RSM.
В этом варианте, помимо символьных вершин, вводятся промежуточные узлы и узлы-диапазоны, соответствующие состояниям и переходам рекурсивного автомата.

Будем придерживаться соглашения о позициях между символами, введённого в определении~\ref{def:PositionInString}: для цепочки длины $n$ позиции нумеруются от $0$ до $n$, и позицией символа считается левая позиция.

\begin{definition}[Сжатое представление леса разбора]
	\label{def:SPPF}
	Пусть $\omega$~--- цепочка длины $n$ над алфавитом $\Sigma$.
    Пусть $G = \langle \Sigma, N, P, S \rangle$~--- КС-грамматика и $\mathcal{R} = \langle \mathcal{N},\Sigma,B,B_S,Q,Q_S \rangle$~--- расширенный рекурсивный автомат для $G$ (раздел~\ref{sec:RecursiveAutomata}).
    \emph{Сжатым представлением леса разбора} (Shared Packed Parse Forest, SPPF) называется ориентированный граф, вершины которого делятся на следующие типы.
	\begin{enumerate}
		\item \emph{Терминальный узел}. Помечен символом $t \in \Sigma$ и парой позиций $i, i+1$, где $i \in [0 \rng n-1]$.
		Обозначается $t_{i,i+1}$ и представляет чтение терминала $t$ между позициями $i$ и $i+1$.
		\item \emph{Нетерминальный узел}. Помечен символом $N \in \mathcal{N}$ и парой позиций $i, j$, где $i, j \in [0 \rng n]$.
		Обозначается $N_{i,j}$ и представляет корень поддерева, соответствующего выводу подцепочки $\omega_{i,j}$ из нетерминала $N$.
		\item \emph{Промежуточный узел}. Помечен промежуточной точкой $I_{q,i}$, где $q \in Q$, $i \in [0 \rng n]$. Соответствует склейке двух диапазонов~--- левого и правого. Всегда имеет ровно двух детей, оба~--- узлы-диапазоны. Нижний индекс --- точка склейки.
		\item \emph{Узел-диапазон}. Помечен диапазоном $R^{q_i,i}_{q_j,j}$, где $q_i,q_j \in Q$, $i,j \in [0 \rng n]$. Обозначает, что алгоритм построил цепочку переходов от конфигурации $(q_i,i,\_)$ до конфигурации $(q_j,j,\_)$. Представляет все возможные способы получить данный диапазон и может иметь произвольное число детей. Дочерними узлами могут быть терминальные, нетерминальные, промежуточные и эпсилон-узлы. Узлы этого типа \emph{переиспользуются}: если один и тот же диапазон может быть получен несколькими способами, все они <<подвешиваются>> к единственному диапазонному узлу.
		\item \emph{Эпсилон-узел}. Обозначается $\varepsilon_i$, где $i \in [0 \rng n]$, и представляет пустое поддерево в позиции $i$, возникающее, когда нетерминал выводит пустую строку.
	\end{enumerate}
\end{definition}

В самом общем виде корнем SPPF для выводимой цепочки длины $n$ является диапазонный узел $R^{q'_0,0}_{q'_1,n}$, где $q'_0$~--- начальное состояние блока для правила $S' \to S \ \$ $, за счёт которого строится расширенная грамматика (и RSM), $q'_1$~--- его финальное состояние.
Листьями являются терминальные узлы и эпсилон-узлы.
\emph{Упаковка альтернатив} реализуется за счёт того, что один диапазонный, нетерминальный или промежуточный узел может иметь несколько дочерних узлов, соответствующих различным способам вывода.

Данное определение использует идеи бинаризованного SPPF: промежуточные узлы как раз соответствуют бинаризации правил грамматики.

Приведём ряд примеров.
Далее будем использовать следующую графическую нотацию: нетерминальные и терминальные узлы~--- прямоугольники со скруглёнными углами,
диапазонные узлы~--- прямоугольники, промежуточные~--- овалы.

\begin{example}
	\label{exmpl:sppf_unambiguous}
	Рассмотрим однозначную грамматику
	\[
	G_0 = \langle \{a,b\}, \{S\}, S, \{S \to a \ S \ b \ S \mid \varepsilon\} \rangle
	\]
	и цепочку $abab$ (позиции $0 \rng 4$).
	Рекурсивный автомат для $G_0$, полученный через UCFS, показан на~рис.~\ref{fig:sppf_rsm_01}.
	Состояние $q_0$ является начальным и допускающим (что соответствует $\varepsilon$-выводу), $q_4$~--- допускающим.
	Переходы $q_1 \xrightarrow{q_0} q_2$ и $q_3 \xrightarrow{q_0} q_4$ соответствуют рекурсивному вызову нетерминала $S$.

	\begin{marginfigure}
		\begin{center}
			\resizebox{\marginparwidth}{!}{\begin{tikzpicture}[->,>=stealth']
    \node[initial,state]   (F)              {$q'_0$};
    \node[state]           (G) [right = of F] {$q'_1$};
    \node[accepting,state] (H) [right = of G] {$q'_2$};
    \node[draw=black, fit= (F) (G) (H), inner sep=0.25cm] (J) {};
    \node[below right] at (J.north west) {$S'$};

    \path (F) edge[below] node {$q_0$} (G)
          (G) edge[below] node {$\$$} (H);

    \node[initial,accepting,state] (A) [below = 2.6cm of F] {$q_0$};
    \node[state]                    (B) [right = of A] {$q_1$};
    \node[state]                    (C) [right = of B] {$q_2$};
    \node[state]                    (D) [right = of C] {$q_3$};
    \node[accepting,state]          (E) [right = of D] {$q_4$};
    \node [draw=black, fit= (A) (B) (C) (D) (E), inner sep=0.25cm] (K) {};
    \node[below right] at (K.north west) {$S$};

    \path (A) edge[below] node {$a$}  (B)
          (B) edge[below] node {$q_0$} (C)
          (C) edge[below] node {$b$}  (D)
          (D) edge[below] node {$q_0$} (E);
\end{tikzpicture}
}
		\end{center}
		\caption{Рекурсивный автомат для грамматики $G_0$: $S \to aSbS \mid \varepsilon$.}
		\label{fig:sppf_rsm_01}
	\end{marginfigure}

	Поскольку грамматика однозначна, для данной цепочки существует ровно одно дерево вывода, и SPPF вырождается в дерево (без разделяемых подграфов).
	SPPF, построенный для данной грамматики и цепочки, представлен на~рис.~\ref{fig:sppf_01}.
	Структура SPPF соответствует дереву вывода: нетерминальные узлы $S_{i,j}$ соединены через диапазонные узлы (прямоугольники) и промежуточные узлы (овалы) с терминальными листьями $a_{i,j}$, $b_{i,j}$ и $\varepsilon$-листьями.

	\begin{figure}[H]
		\begin{center}
			\resizebox{0.7\textwidth}{!}{\begin{tikzpicture}[shorten >=1pt,node distance=0.6cm and 1cm]
    \node[symbol_node] (us2us0) {$S_{0,4}$};
    \node[prod_node] (us2us1) [below=of us2us0] {$R^{q_0,0}_{q_4,4}$};
    \node[intermediate_node] (us2us12) [below=of us2us1] {$I_{q_3,2}$};
    \node[prod_node] (us2us23) [below left=of us2us12] {$R^{q_0,0}_{q_3,2}$};
    \node[prod_node] (us2us29) [right=of us2us23] {$R^{q_3,2}_{q_4,4}$};
    \node[intermediate_node] (us2us30) [below=of us2us23] {$I_{q_2,1}$};
    \node[symbol_node] (us2us31) [right=of us2us30] {$S_{2,4}$};
    \node[prod_node] (us2us32) [below left=of us2us30] {$R^{q_0,0}_{q_2,1}$};
    \node[prod_node] (us2us33) [right=of us2us32] {$R^{q_2,1}_{q_3,2}$};
    \node[prod_node] (us2us34) [right=of us2us33] {$R^{q_0,2}_{q_4,4}$};
    \node[intermediate_node] (us2us2) [below=of us2us32] {$I_{q_1,1}$};
    \node[symbol_node] (us2us3) [right=of us2us2] {$b_{1,2}$};
    \node[intermediate_node] (us2us4) [right=of us2us3] {$I_{q_3,4}$};
    \node[prod_node] (us2us5) [below left=of us2us2] {$R^{q_0,0}_{q_1,1}$};
    \node[prod_node] (us2us6) [right=of us2us5] {$R^{q_1,1}_{q_2,1}$};
    \node[prod_node] (us2us7) [right=of us2us6] {$R^{q_0,2}_{q_3,4}$};
    \node[prod_node] (us2us8) [right=of us2us7] {$R^{q_3,4}_{q_4,4}$};
    \node[symbol_node] (us2us9) [below=of us2us5] {$a_{0,1}$};
    \node[symbol_node] (us2us10) [right=of us2us9] {$S_{1,1}$};
    \node[intermediate_node] (us2us11) [right=of us2us10] {$I_{q_2,3}$};
    \node[symbol_node] (us2us13) [right=of us2us11] {$S_{4,4}$};
    \node[prod_node] (us2us14) [below=of us2us10] {$R^{q_0,1}_{q_0,1}$};
    \node[prod_node] (us2us15) [right=of us2us14] {$R^{q_0,2}_{q_2,3}$};
    \node[prod_node] (us2us16) [right=of us2us15] {$R^{q_2,3}_{q_3,4}$};
    \node[prod_node] (us2us17) [right=of us2us16] {$R^{q_0,4}_{q_0,4}$};
    \node[symbol_node] (us2us18) [below=of us2us14] {$\varepsilon_{1}$};
    \node[intermediate_node] (us2us19) [right=of us2us18] {$I_{q_1,3}$};
    \node[symbol_node] (us2us20) [right=of us2us19] {$b_{3,4}$};
    \node[symbol_node] (us2us21) [right=of us2us20] {$\varepsilon_{4}$};
    \node[prod_node] (us2us22) [below right=of us2us18] {$R^{q_0,2}_{q_1,3}$};
    \node[prod_node] (us2us24) [right=of us2us22] {$R^{q_1,3}_{q_2,3}$};
    \node[symbol_node] (us2us25) [below=of us2us22] {$a_{2,3}$};
    \node[symbol_node] (us2us26) [right=of us2us25] {$S_{3,3}$};
    \node[prod_node] (us2us27) [below=of us2us26] {$R^{q_0,3}_{q_0,3}$};
    \node[symbol_node] (us2us28) [below=of us2us27] {$\varepsilon_{3}$};

    \path[->]
        (us2us0) edge (us2us1)
        (us2us1) edge (us2us12)
        (us2us2) edge (us2us5)
        (us2us2) edge (us2us6)
        (us2us4) edge (us2us7)
        (us2us4) edge (us2us8)
        (us2us5) edge (us2us9)
        (us2us6) edge (us2us10)
        (us2us7) edge (us2us11)
        (us2us8) edge (us2us13)
        (us2us10) edge (us2us14)
        (us2us11) edge (us2us15)
        (us2us11) edge (us2us16)
        (us2us12) edge (us2us23)
        (us2us12) edge (us2us29)
        (us2us13) edge (us2us17)
        (us2us14) edge (us2us18)
        (us2us15) edge (us2us19)
        (us2us16) edge (us2us20)
        (us2us17) edge (us2us21)
        (us2us19) edge (us2us22)
        (us2us19) edge (us2us24)
        (us2us22) edge (us2us25)
        (us2us23) edge (us2us30)
        (us2us24) edge (us2us26)
        (us2us26) edge (us2us27)
        (us2us27) edge (us2us28)
        (us2us29) edge (us2us31)
        (us2us30) edge (us2us32)
        (us2us30) edge (us2us33)
        (us2us31) edge (us2us34)
        (us2us32) edge (us2us2)
        (us2us33) edge (us2us3)
        (us2us34) edge (us2us4)
    ;
\end{tikzpicture}
}
		\end{center}
		\caption{SPPF для цепочки $abab$ в однозначной грамматике $G_0$.}
		\label{fig:sppf_01}
	\end{figure}
\end{example}

\begin{example}
	\label{exmpl:sppf_ambiguous}
	Рассмотрим неоднозначную грамматику
	\[
	G_1 = \langle \{a,b\}, \{S\}, S, \{S \to S \ S \mid a \ S \ b \mid \varepsilon\} \rangle
	\]
	и цепочку $ab$ (позиции $0 \rng 2$).
	Рекурсивный автомат для $G_1$ показан на~рис.~\ref{fig:sppf_rsm_02}.
	Из начального состояния $q_0$ (допускающего $\varepsilon$) исходят два перехода: $q_0 \xrightarrow{q_0} q_1$ (первый символ альтернативы $SS$) и $q_0 \xrightarrow{a} q_2$ (начало альтернативы $aSb$).
	Обе альтернативы сходятся в допускающем состоянии $q_4$.

	\begin{marginfigure}
		\begin{center}
			\resizebox{\marginparwidth}{!}{\begin{tikzpicture}[->,>=stealth']
    \node[initial,state]   (F)              {$q'_0$};
    \node[state]           (G) [right = of F] {$q'_1$};
    \node[accepting,state] (H) [right = of G] {$q'_2$};
    \node[draw=black, fit= (F) (G) (H), inner sep=0.25cm] (J) {};
    \node[below right] at (J.north west) {$S'$};

    \path (F) edge[below] node {$q_0$} (G)
          (G) edge[below] node {$\$$} (H);

    \node[initial,accepting,state] (A) [below = 2.6cm of F] {$q_0$};
    \node[state]                    (B) [above right = 0.3cm and 1.6cm of A] {$q_1$};
    \node[state]                    (C) [below right = 0.3cm and 1.6cm of A] {$q_2$};
    \node[state]                    (D) [right = of C]  {$q_3$};
    \node[accepting,state]          (E) [right = of B]  {$q_4$};
    \node [draw=black, fit= (A) (B) (C) (D) (E), inner sep=0.25cm] (K) {};
    \node[below right] at (K.north west) {$S$};

    \path (A) edge[left]  node[pos=0.35] {$q_0$} (B)
          (B) edge[below] node {$q_0$} (E)
          (A) edge[left]  node[pos=0.55] {$a$}   (C)
          (C) edge[below] node {$q_0$} (D)
          (D) edge[right] node {$b$}   (E);
\end{tikzpicture}
}
		\end{center}
		\caption{Рекурсивный автомат для грамматики $G_1$: $S \to SS \mid aSb \mid \varepsilon$.}
		\label{fig:sppf_rsm_02}
	\end{marginfigure}

	Цепочка $ab$ в данной грамматике имеет бесконечное число выводов за счёт $S \derives[+] S$. Таким конструкциям как раз будут соответствовать циклы в SPPF.

	SPPF для $G_1$ и цепочки $ab$ показан на~рис.~\ref{fig:sppf_02}.
	Здесь проявляются две ключевые особенности SPPF.
	\begin{enumerate}
		\item \emph{Переиспользование поддеревьев.} Узел $S_{1,1}$ (соответствующий $\varepsilon$-выводу) используется в обоих деревьях: и при разборе через $aSb$, и через $SS$.
		\item \emph{Циклы для $\varepsilon$.} Узлы $S_{0,0}$, $S_{1,1}$, $S_{2,2}$ соответствуют $\varepsilon$-выводам. Каждый из них может быть получен через правило $S \to SS$ как $S_{i,i} \to S_{i,i} \cdot S_{i,i}$, что порождает цикл в SPPF.
	\end{enumerate}
	Заметим, что применение правила $S \to SS$ к корневому узлу $S_{0,2}$ приводит к тому, что среди его дочерних узлов снова появляется $S_{0,2}$: первое разбиение~--- $S_{0,0} \cdot S_{0,2}$, второе~--- $S_{0,2} \cdot S_{2,2}$.
	Циклы означают, что конечный SPPF-граф представляет бесконечное семейство выводов: любое число дополнительных витков рекурсии через $S \to SS$ с $\varepsilon$-поддеревьями порождает формально различные, но эквивалентные выводы.

	\begin{figure}[H]
		\begin{center}
			\resizebox{0.95\textwidth}{!}{\begin{tikzpicture}[shorten >=1pt,node distance=0.6cm and 1cm]
    \node[symbol_node] (us1us0) {$S_{0,2}$};
    \node[prod_node] (us1us1) [below=of us1us0] {$R^{q_0,0}_{q_4,2}$};
    \node[intermediate_node] (us1us12) [below left=of us1us1] {$I_{q_3,1}$};
    \node[intermediate_node] (us1us23) [right=of us1us12] {$I_{q_1,0}$};
    \node[intermediate_node] (us1us34) [right=of us1us23] {$I_{q_1,2}$};
    \node[prod_node] (us1us36) [below left=of us1us12] {$R^{q_0,0}_{q_3,1}$};
    \node[prod_node] (us1us37) [right=of us1us36] {$R^{q_3,1}_{q_4,2}$};
    \node[prod_node] (us1us38) [right=of us1us37] {$R^{q_0,0}_{q_1,0}$};
    \node[prod_node] (us1us39) [right=of us1us38] {$R^{q_1,0}_{q_4,2}$};
    \node[prod_node] (us1us2) [right=of us1us39] {$R^{q_1,2}_{q_4,2}$};
    \node[prod_node] (us1us40) [right=of us1us2] {$R^{q_0,0}_{q_1,2}$};
    \node[intermediate_node] (us1us3) [below=of us1us36] {$I_{q_2,1}$};
    \node[symbol_node] (us1us4) [right=of us1us3] {$b_{1,2}$};
    \node[symbol_node] (us1us5) [right=of us1us4] {$S_{0,0}$};
    \node[symbol_node] (us1us6) [right=of us1us5] {$S_{0,2}$};
    \node[symbol_node] (us1us8) [right=of us1us6] {$S_{2,2}$};
    \node[symbol_node] (us1us7) [right=of us1us8] {$S_{0,2}$};
    \node[prod_node] (us1us9) [below left=of us1us3] {$R^{q_0,0}_{q_2,1}$};
    \node[prod_node] (us1us10) [right=of us1us9] {$R^{q_2,1}_{q_3,1}$};
    \node[prod_node] (us1us13) [right=of us1us10] {$R^{q_0,0}_{q_4,0}$};
    \node[prod_node] (us1us15) [right=of us1us13] {$R^{q_0,2}_{q_4,2}$};
    \node[symbol_node] (us1us16) [below=of us1us9] {$a_{0,1}$};
    \node[symbol_node] (us1us17) [right=of us1us16] {$S_{1,1}$};
    \node[intermediate_node] (us1us19) [right=of us1us17] {$I_{q_1,0}$};
    \node[intermediate_node] (us1us21) [right=of us1us19] {$I_{q_1,2}$};
    \node[prod_node] (us1us24) [below=of us1us17] {$R^{q_0,1}_{q_4,1}$};
    \node[prod_node] (us1us25) [right=of us1us24] {$R^{q_1,0}_{q_4,0}$};
    \node[prod_node] (us1us26) [right=of us1us25] {$R^{q_0,2}_{q_1,2}$};
    \node[intermediate_node] (us1us28) [below=of us1us24] {$I_{q_1,1}$};
    \node[symbol_node] (us1us29) [right=of us1us28] {$S_{0,0}$};
    \node[symbol_node] (us1us30) [right=of us1us29] {$S_{2,2}$};
    \node[prod_node] (us1us31) [below left=of us1us28] {$R^{q_0,1}_{q_1,1}$};
    \node[prod_node] (us1us32) [right=of us1us31] {$R^{q_1,1}_{q_4,1}$};
    \node[prod_node] (us1us11) [right=of us1us32] {$R^{q_0,0}_{q_0,0}$};
    \node[prod_node] (us1us14) [right=of us1us11] {$R^{q_0,2}_{q_0,2}$};
    \node[symbol_node] (us1us33) [below=of us1us31] {$S_{1,1}$};
    \node[symbol_node] (us1us35) [right=of us1us33] {$S_{1,1}$};
    \node[symbol_node] (us1us18) [right=of us1us35] {$\varepsilon_{0}$};
    \node[symbol_node] (us1us20) [right=of us1us18] {$\varepsilon_{2}$};
    \node[prod_node] (us1us22) [below=of us1us35] {$R^{q_0,1}_{q_0,1}$};
    \node[symbol_node] (us1us27) [below=of us1us22] {$\varepsilon_{1}$};

    \path[->]
        (us1us0) edge (us1us1)
        (us1us1) edge (us1us12)
        (us1us1) edge (us1us23)
        (us1us1) edge (us1us34)
        (us1us2) edge (us1us8)
        (us1us3) edge (us1us9)
        (us1us3) edge (us1us10)
        (us1us5) edge [bend right] (us1us11)
        (us1us5) edge (us1us13)
        (us1us6) edge [bend right] (us1us1)
        (us1us7) edge [bend right] (us1us1)
        (us1us8) edge [bend right] (us1us14)
        (us1us8) edge (us1us15)
        (us1us9) edge (us1us16)
        (us1us10) edge (us1us17)
        (us1us11) edge (us1us18)
        (us1us12) edge (us1us36)
        (us1us12) edge (us1us37)
        (us1us13) edge (us1us19)
        (us1us14) edge (us1us20)
        (us1us15) edge (us1us21)
        (us1us17) edge [bend right] (us1us22)
        (us1us17) edge (us1us24)
        (us1us19) edge (us1us25)
        (us1us19) edge [bend right] (us1us38)
        (us1us21) edge [bend right] (us1us2)
        (us1us21) edge (us1us26)
        (us1us22) edge (us1us27)
        (us1us23) edge (us1us38)
        (us1us23) edge (us1us39)
        (us1us24) edge (us1us28)
        (us1us25) edge (us1us29)
        (us1us26) edge (us1us30)
        (us1us28) edge (us1us31)
        (us1us28) edge (us1us32)
        (us1us29) edge (us1us11)
        (us1us29) edge [bend right] (us1us13)
        (us1us30) edge (us1us14)
        (us1us30) edge [bend right] (us1us15)
        (us1us31) edge (us1us33)
        (us1us32) edge (us1us35)
        (us1us33) edge (us1us22)
        (us1us33) edge [bend right] (us1us24)
        (us1us34) edge (us1us2)
        (us1us34) edge (us1us40)
        (us1us35) edge (us1us22)
        (us1us35) edge [bend right] (us1us24)
        (us1us36) edge (us1us3)
        (us1us37) edge (us1us4)
        (us1us38) edge (us1us5)
        (us1us39) edge (us1us6)
        (us1us40) edge (us1us7)
    ;
\end{tikzpicture}
}
		\end{center}
		\caption{SPPF для цепочки $ab$ в неоднозначной грамматике $G_1$. Показано переиспользование $\varepsilon$-узлов и циклы через правило $S \to SS$.}
		\label{fig:sppf_02}
	\end{figure}
\end{example}

\begin{example}
	\label{exmpl:sppf_ebnf}
	Рассмотрим грамматику в EBNF-форме $S \to (a \ S \ b)^*$, задающую то же семейство языков, и цепочку $abab$ (позиции $0 \rng 4$).
	Рекурсивный автомат для неё показан на~рис.~\ref{fig:sppf_rsm_03}.
	Здесь $q_0$~--- начальное и допускающее состояние (базовый случай~--- пустая строка), а цикл $q_0 \xrightarrow{a} q_1 \xrightarrow{q_0} q_2 \xrightarrow{b} q_0$ соответствует одной итерации $(a S b)$.

	\begin{marginfigure}
		\begin{center}
			\resizebox{\marginparwidth}{!}{\input{part_02_Foundations/chapter_06_ContextFreeLanguages/figures/05_SPPF/05_SPPF_rsm_03.tex}}
		\end{center}
		\caption{Рекурсивный автомат для грамматики $S \to (a S b)^*$.}
		\label{fig:sppf_rsm_03}
	\end{marginfigure}

	SPPF для данной грамматики и цепочки представлен на~рис.~\ref{fig:sppf_03}.
	Как и в случае дерева вывода для EBNF-грамматики (раздел~\ref{sec:DerivTree}), если удалить вспомогательные диапазонные и промежуточные узлы, оставив только связь <<родительский нетерминал~--- его дочерние нетерминалы и терминалы>>, то последовательность меток дочерних узлов будет образовывать цепочку, принимаемую блоком рекурсивного автомата для данного нетерминала.

	\begin{figure}[H]
		\begin{center}
			\resizebox{0.7\textwidth}{!}{\begin{tikzpicture}[shorten >=1pt,node distance=1.2cm]
    \node[symbol_node] (us2us0) {$S_{0,4}$};
    \node[prod_node] (us2us1) [below=of us2us0] {$R^{q_0,0}_{q_0,4}$};
    \node[intermediate_node] (us2us12) [below=of us2us1] {$I_{q_2,3}$};
    \node[prod_node] (us2us20) [below left=of us2us12] {$R^{q_0,0}_{q_2,3}$};
    \node[prod_node] (us2us21) [right=of us2us20] {$R^{q_2,3}_{q_0,4}$};
    \node[intermediate_node] (us2us22) [below=of us2us20] {$I_{q_1,3}$};
    \node[symbol_node] (us2us23) [right=of us2us22] {$b_{3,4}$};
    \node[prod_node] (us2us24) [below left=of us2us22] {$R^{q_0,0}_{q_1,3}$};
    \node[prod_node] (us2us25) [right=of us2us24] {$R^{q_1,3}_{q_2,3}$};
    \node[intermediate_node] (us2us26) [below=of us2us24] {$I_{q_0,2}$};
    \node[symbol_node] (us2us2) [right=of us2us26] {$S_{3,3}$};
    \node[prod_node] (us2us3) [below left=of us2us26] {$R^{q_0,0}_{q_0,2}$};
    \node[prod_node] (us2us4) [right=of us2us3] {$R^{q_0,2}_{q_1,3}$};
    \node[prod_node] (us2us5) [right=of us2us4] {$R^{q_0,3}_{q_0,3}$};
    \node[intermediate_node] (us2us6) [below=of us2us3] {$I_{q_2,1}$};
    \node[symbol_node] (us2us7) [right=of us2us6] {$a_{2,3}$};
    \node[symbol_node] (us2us8) [right=of us2us7] {$\varepsilon_{3}$};
    \node[prod_node] (us2us9) [below left=of us2us6] {$R^{q_0,0}_{q_2,1}$};
    \node[prod_node] (us2us10) [right=of us2us9] {$R^{q_2,1}_{q_0,2}$};
    \node[intermediate_node] (us2us11) [below=of us2us9] {$I_{q_1,1}$};
    \node[symbol_node] (us2us13) [right=of us2us11] {$b_{1,2}$};
    \node[prod_node] (us2us14) [below left=of us2us11] {$R^{q_0,0}_{q_1,1}$};
    \node[prod_node] (us2us15) [right=of us2us14] {$R^{q_1,1}_{q_2,1}$};
    \node[symbol_node] (us2us16) [below=of us2us14] {$a_{0,1}$};
    \node[symbol_node] (us2us17) [right=of us2us16] {$S_{1,1}$};
    \node[prod_node] (us2us18) [below=of us2us17] {$R^{q_0,1}_{q_0,1}$};
    \node[symbol_node] (us2us19) [below=of us2us18] {$\varepsilon_{1}$};

    \path[->]
        (us2us0) edge (us2us1)
        (us2us1) edge (us2us12)
        (us2us2) edge (us2us5)
        (us2us3) edge (us2us6)
        (us2us4) edge (us2us7)
        (us2us5) edge (us2us8)
        (us2us6) edge (us2us9)
        (us2us6) edge (us2us10)
        (us2us9) edge (us2us11)
        (us2us10) edge (us2us13)
        (us2us11) edge (us2us14)
        (us2us11) edge (us2us15)
        (us2us12) edge (us2us20)
        (us2us12) edge (us2us21)
        (us2us14) edge (us2us16)
        (us2us15) edge (us2us17)
        (us2us17) edge (us2us18)
        (us2us18) edge (us2us19)
        (us2us20) edge (us2us22)
        (us2us21) edge (us2us23)
        (us2us22) edge (us2us24)
        (us2us22) edge (us2us25)
        (us2us24) edge (us2us26)
        (us2us25) edge (us2us2)
        (us2us26) edge (us2us3)
        (us2us26) edge (us2us4)
    ;
\end{tikzpicture}
}
		\end{center}
		\caption{SPPF для цепочки $abab$ в грамматике $S \to (a S b)^*$. Промежуточные узлы $I_{S_0}$, $I_{S_1}$, $I_{S_2}$ соответствуют состояниям рекурсивного автомата.}
		\label{fig:sppf_03}
	\end{figure}
\end{example}

Применение SPPF для представления результатов КС-запросов к графам обсуждается в главе~\ref{chpt:CFPQ}, в частности обобщение определения~\ref{def:SPPF} на графовый вход (с заменой позиций на вершины графа).
